% !PRINTED 424
\weberpage{424}
%<<CONT
\weq{3}{\Sigma N' = 2\sum^{\delta}\sum_{\partial' > a'}^{\partial'} \frac{\partial'}{e'}\varphi(e') + \sum^{n'} \varphi(\sqrt{n'}),}
if $\delta^2$ runs through all the square factors of $n$ (with the
possible exception of $n$ itself).

If we now set
\[
\partial'\delta = \partial, \quad a'\delta = a, \quad e'\delta = e,
\]
then $a\partial = n$, and $e$ is the greatest common divisor of $a$ and
$\partial$; at the same time $\partial' > a'$ has as its consequence
that $\partial > a$. Conversely, from every decomposition $a\partial$
of $n$ and from every divisor $e'$ of $e$ one obtains a decomposition
$a'\partial'$ of $n'$, where
\[
\delta = \frac{e}{e'}, \quad \partial' = \frac{\partial e'}{e}, \quad a' = \frac{a e'}{e}
\]
holds. Accordingly
\weq{4}{\Sigma N' = 2\sum_{\partial>a}^{\partial} \frac{\partial}{e} \sum \varphi(e') + \sum^{n'} \varphi(\sqrt{n'}).}

In this we now make use of the number-theoretic relation:
$\sum \varphi(\partial) = n$, in which the sum refers to all divisors
of $n$. If we first assume that $n$ is not a square, then none of the
$n'$ are squares either, and $\varphi(\sqrt{n'})=0$. Hence, since
$\sum \varphi(e') = e$, we obtain
\weq{5}{\Sigma N' = 2\sum_{\partial>\sqrt{n}} \partial.}

If on the other hand $n$ is a square, then $e'$ still runs through all
the divisors of $e$, but $\sqrt{n'}$ runs through all the divisors of
$\sqrt{n}$, with the exception of 1; from this there follows for this
case:
\weq{6}{\Sigma N' = 2\sum_{\partial>\sqrt{n}} \partial + \sqrt{n} - 1.}

In both cases $\partial$ runs through all the divisors of $n$ that are
\gesperrt{greater} than $\sqrt{n}$.

We have now still to form the sum of the expressions $\Sigma k h$ for
the various values of $n'$ or $\delta$. But if we set
$n = \delta^2 n'$, then from every proper solution of
\weq{7}{4n' = y'^2 + m x'^2}
we obtain, by multiplication by $\delta^2$, a solution (improper, if
$\delta > 1$) of
\weq{8}{4n = y^2 + m x^2,}
namely
\[
y = \delta y', \quad x = \delta x'.
\]
